\documentclass[t]{beamer}

\usepackage[utf8]{inputenc}
\usepackage[T1]{fontenc}
\usepackage{polski}
\usepackage{graphicx} 
\usepackage{epstopdf}
\usepackage{listings}

\definecolor{mint}{rgb}{0.96, 1, 0.98}
\lstset
{
    language=[LaTeX]TeX,
    inputencoding=utf8,
    breaklines=true,
    basicstyle=\tt\scriptsize,
    keywordstyle=\color{red},
    identifierstyle=\color{blue},
    frame=lines,
    backgroundcolor=\color{mint}
}

%%%Komendy
%%%Beamer
\newcommand{\beamer}{
{\rmfamily\textsc{Beamer}}
}

\usetheme{Copenhagen}
\usecolortheme{crane}
\setbeamertemplate{navigation symbols}{} 
\expandafter\def\expandafter\insertshorttitle\expandafter{%
  \insertshorttitle\hfill%
  \insertframenumber\,/\,\inserttotalframenumber}
\mode<presentation>
\title{Seminarium dyplomowe inżynierskie}
\subtitle{Temat: Aplikacja WYSIWYG do tworzenia prezentacji multimedialnych z wykorzystaniem klasy dokumentów \beamer \LaTeX-a}
\author[J.A. Kumor]{Jan Adam Kumor}
\institute[IAiIS]{\begin{minipage}[c]{2.0cm}
			\includegraphics[width = 1.5cm]{IAiISLogo.eps}
		\end{minipage} Instytut Automatyki i Informatyki Stosowanej\\[0.6cm]}
\date[2015-10-23]{23 października 2015}

\usepackage{graphicx}

\begin{document}

\AtBeginSubsection{
	\begin{frame}
		\frametitle{Konspekt}
		\tableofcontents[currentsection, currentsubsection]
	\end{frame}
}

\begin{frame}
	\titlepage
\end{frame}

\section{Wprowadzenie}
\subsection{}

\begin{frame}[t]
	\frametitle{Kto, co i gdzie?}
	\begin{description}[style=multiline]
		\item[Autor:] Jan Kumor, semestr 7. specjalność SID
		\item[Opiekun:] prof. nzw. dr hab. inż. Maciej Ławryńczuk \\ Zespół Technik Sterowania \\ Zakład Automatyki i Inżynierii Oprogramowania \\ Instytut Automatyki i Informatyki Stosowanej
		\item[Temat:] \emph{Aplikacja WYSIWYG do tworzenia prezentacji multimedialnych z wykorzystaniem klasy dokumentów} \beamer{} \LaTeX-a
	\end{description}
\end{frame}

\section{Omówienie tematyki pracy}

\subsection{Podstawowe terminy}

\begin{frame}
	\frametitle{Systemy WYSIWYG}
	\begin{block}{\emph{WYSIWYG}}
		skrót od ang. \emph{What You See Is What You Get}
	\end{block}
	\begin{itemize}
		\item Akronim ten określa metodę pracy edytorów
		\item System WYSIWYG wyświetla zawartość w trakcie edycji w sposób odpowiadający efektowi końcowemu
		\item Rodzaje systemów WYSIWYG:
			\begin{itemize}
				\item procesory tekstu np. MS Word
				\item edytory grafiki np. GIMP
				\item programy CAD np. Autodesk AutoCAD
				\item edytory stron www np. tryb graficzny w Pajączek
			\end{itemize}
	\end{itemize}
\end{frame}
\begin{frame}
	\frametitle{\TeX\ i \LaTeX\ }
	\begin{quotation}
		\TeX\ – komputerowy system profesjonalnego składu drukarskiego, obejmujący zarówno specjalny język, jak i kompilator przygotowujący pliki w formatach wymaganych przez urządzenia graficzne (drukarki, naświetlarki).\footnote{Wikipedia, https://pl.wikipedia.org/wiki/TeX}
	\end{quotation}
	\begin{quotation}
	[...] LaTeX nie jest samodzielnym środowiskiem programistycznym. Jest to jedynie zestaw makr stanowiących nadbudowę dla systemu składu TeX, automatyzujących wiele czynności związanych z procesem poprawnego składania tekstu.\footnote{Wikipedia, https://pl.wikipedia.org/wiki/LaTeX}
	\end{quotation}
\end{frame}

\begin{frame}[fragile]
	\frametitle{Fragment przykładowego dokumentu \LaTeX}
	\begin{lstlisting}
		\section{Szablony predefiniowanych ukladow slajdow}

\subsection{Uwagi ogolne}
W rozdziale tym zaprezentowane zostaly szablony slajdow, ktore bedzie mogl wybrac uzytkownik w aplikacji \projektor. Opis kazdego szablonu sklada sie z:
\begin{itemize}
        \item{wizualizacji slajdu}
        \item{krotkiego opisu tekstowego}
        \item{parametrow (opcji) za pomoca, ktorych uzytkownik bedzie mial mozliwosc spersonalizowania slajdu}
        \item{kodu \LaTeX, z ktorego generowany jest slajd}
\end{itemize}
	\end{lstlisting}
\end{frame}

\begin{frame}[plain]
	\begin{figure}[T]
		\flushleft
		\hspace{-1cm}\includegraphics[width=\textwidth]{LaTeX_diagram.png}
		\caption{Relacje między źródłami \TeX\ i \LaTeX\ a formatami wyjściowymi. Źródło: Wikibooks. Latex — wikibooks, the free textbook project, \url{https://en.wikibooks.org/wiki/LaTeX/Basics\#Compilation}}
		\label{}
	\end{figure}
\end{frame}

\begin{frame}
	\frametitle{Klasa dokumentów \beamer\ }
	\begin{itemize}
		\item Dedykowana tworzenia prezentacji multimedialnych
		\item Dostarcza środowisko \texttt{frame} (z ang. \emph{frame}  - ramka)
		\item Środowisko \texttt{frame} dostarcza mechanizm \texttt{overlay specifications} (z ang. \emph{overlay} - nakładka)
		\begin{alertblock}{Uwaga}
			Ramka nie jest tożsama ze slajdem prezentacji. Z jednej ramki może zostać wygenerowane kilka slajdów.
		\end{alertblock}
		\item Możliwość wygenerowania tzw. \emph{handouts} czyli slajdów z pominięciem \texttt{overlay specifications} 
		\item Pakiet \emph{beamerarticle} pozwala na wygenerowanie treści prezentacji w formie tekstowej
	\end{itemize}
\end{frame}

\begin{frame}[fragile]
	\frametitle{Środowisko \texttt{frame} i mechanizm \texttt{overlay specifications}}
	\begin{columns}[T]
		\begin{column}{.55\textwidth}
			\begin{lstlisting}
				\begin{frame}
					\frametitle{Przyklad overlay specifications}
					\begin{itemize}
						\item <1-> \alert<1>{Punkt pierwszy}
						\item <2-> \alert<2>{Punkt drugi}
					\end{itemize}
				\end{frame }
			\end{lstlisting}
		\end{column}
		\begin{column}{.4\textwidth}
			\framebox{
				\includegraphics[page=1,width=.9\textwidth]{overlay_example.pdf}
			}\\
			\framebox{
				\includegraphics[page=2,width=.9\textwidth]{overlay_example.pdf}
			}
		\end{column}		
	\end{columns}
\end{frame}

\subsection{Istniejące rozwiązania}

\begin{frame}[T]
	\frametitle{Przykład 1. Gummi \footnote{\url{http://gummi.midnightcoding.org/}}}
		\begin{columns}[T]
			\begin{column}{.15\textwidth}
				\includegraphics[width=2cm]{gummi_logo.png}
			\end{column}
			\begin{column}{.85\textwidth}
				\begin{itemize}
					\item Typowe dla dokumentów \LaTeX\ workflow - plik \emph{.tex} i podgląd wyniku kompilacji
					\item Automatyczna kompilacja dokumentu po każdej wprowadzonej zmianie
					\item Dostępne  podstawowe operacje na tekście:
					\begin{itemize}
						\item \textbf{pogrubienie}, \emph{kursywa} oraz \underline{podkreślenie} tekstu
						\item wyrównanie tekstu do lewej i prawej strony oraz wyśrodkowanie
					\end{itemize}
					\item Wsparcie GUI przy wstawianiu grafiki, tabel, macierzy i bibliografii
				\end{itemize}
			\end{column}		
		\end{columns}
\end{frame}

\begin{frame}[T]
	\frametitle{Przykład 1. Gummi}
	\begin{figure}[p]
   	    \centering
		\includegraphics[width=.95\textwidth]{gummi_prtscn.png}
   	    \caption{Okno programu Gummi}
	\end{figure}
\end{frame}


\begin{frame}[t]
	\frametitle{Przykład 2. LyX – The Document Processor \footnote{\url{http://www.lyx.org/}}}
		\begin{columns}[T]
			\begin{column}{.15\textwidth}
				\includegraphics[width=2cm]{lyx_logo.png}
			\end{column}
			\begin{column}{.85\textwidth}
				\begin{itemize}
					\item Wsparcie GUI dla większości funkcji \LaTeX\ 
					\item Wiele szablonów dla poszczególnych klas dokumentów
					\item Teoretycznie nie wymaga znajomości \LaTeX\ 
					\item Wyświetla w przybliżeniu wygląd treści dokumentu docelowego
					\item Zapis plików w formacie tekstowym wzorowanym na xml
					\item Współpraca z wieloma konwerterami
				\end{itemize}
			\end{column}		
		\end{columns}
\end{frame}

\begin{frame}[t]
	\frametitle{Przykład 2. LyX – The Document Processor}
		\begin{figure}[p]
     	    \centering
			\includegraphics[width=.95\textwidth]{lyx_prtscn.png}
     	    \caption{Okno programu LyX}
		\end{figure}
		
\end{frame}

\begin{frame}[t]
	\frametitle{Przykład 3. BaKoMa \TeX\ (True WYSIWYG \LaTeX\ System)\footnote{\url{http://www.bakoma-tex.com/}}}
		\begin{columns}[T]
			\begin{column}{.15\textwidth}
				\includegraphics[width=2cm]{bakoma_logo.png}
			\end{column}
			\begin{column}{.85\textwidth}
				\begin{itemize}
					\item Pakiet aplikacji komercyjnych 
					\item Wsparcie GUI dla większości funckji \LaTeX\ 
					\item Obsługa wielu klas dokumentów z poziomu GUI
					\item Prawdziwy WYSIWYG w czasie rzeczywistym
					\item W tle działa silnik \TeX - tylko dla aktualnie zmienianej strony
					\item Dwukierunkowa synchronizacja między plikiem żródłowym a widokiem
				\end{itemize}
			\end{column}		
		\end{columns}
\end{frame}

\begin{frame}[t]
	\frametitle{Przykład 3. BaKoMa \TeX\ (True WYSIWYG \LaTeX\ System)}
		\begin{figure}[p]
     	    \centering
			\includegraphics[width=.95\textwidth]{bakoma_prtscn.png}
     	    \caption{Okno programu BaKoMa Word}
		\end{figure}
		
\end{frame}

\begin{frame}
	\frametitle{Istniejące rozwiązania - ogólny obraz}
	\begin{itemize}
		\item Istnieje wiele narzędzi wspomagających pracę z \LaTeX-em
		\item O różnym stopniu złożoności i funkcjonalności
		\item Brak rozwiązań dedykowanych stricte edycji dokumentów klasy \beamer		
		\item Jednak większa część funkcji to czysty \LaTeX
	\end{itemize}
	\begin{block}{Wniosek}
		Tworzona aplikacja powinna dostarczać wsparcia dla zestawu funkcji \LaTeX-a wzorowanego na istniejących rozwiązaniach i rozszerzyć je o mechanizmy dedykowane klasie \beamer.
	\end{block}
\end{frame}

\subsection{Sformułowanie zadania}

\begin{frame}
	\frametitle{Dekompozycja problemu}
	W celu łatwiejszej analizy dokonano podziału aplikacji na obszary zgodnie ze schematem.
	\begin{figure}[htp]
		\centering
		\includegraphics[width=.85\textwidth]{obszary.png}
		\label{obsz}
	\end{figure}
\end{frame}

\begin{frame}
	\frametitle{Interfejs do tworzenia prezentacji}
	\begin{itemize}
		\item Typowe funkcje edycyji tekstu i akapitów
		\item Wsparcie przy tworzeniu elementów osadzonych (grafiki, tabele, bloki)
		\item Tworzenie slajdów z użyciem predefiniowanych szablonów
		\item Wsparcie dla mechanizmu \texttt{overlay specifications}
		\item Zarządzanie strukturą dokumentu
		\begin{itemize}
			\item porządkowanie slajdów
			\item tworzenie sekcji i podsekcji
			\item mechanizmy \texttt{AtBeginSection} i \texttt{AtBeginSubsection}
			\item zarządzanie bibliografią
		\end{itemize}
	\end{itemize}
\end{frame}

\begin{frame}[T]
	\frametitle{Przykłady predefiniowanych szablonów}
	\only<1>{
		\begin{columns}[T]
		\begin{column}{.5\textwidth}
			\centering
			\framebox{
				\includegraphics[page=7, width=\textwidth]{templates.pdf}
			}
		\end{column}
		\begin{column}{.5\textwidth}
			\centering
			\framebox{
				\includegraphics[page=14, width=\textwidth]{templates.pdf}
			}
		\end{column}
		\end{columns}
	}
	\only<2>{
			\begin{columns}[T]
		\begin{column}{.5\textwidth}
			\centering
			\framebox{
				\includegraphics[page=22, width=\textwidth]{templates.pdf}
			}
		\end{column}
		\begin{column}{.5\textwidth}
			\centering
			\framebox{
				\includegraphics[page=30, width=\textwidth]{templates.pdf}
			}
		\end{column}
		\end{columns}
	}
\end{frame}

\begin{frame}
	\frametitle{Generator kodu \LaTeX}
	\begin{itemize}
		\item Przekształcenie wewnętreznego modelu na plik .tex
		\item Powiązanie wynikowego pliku z modelem w celu synchronizacji
		\item Ważna wydajność - potencjalnie najczęśćiej wołany fragment kodu własnego
	\end{itemize}
\end{frame}

\begin{frame}
	\frametitle{Moduł prezentacji - wariant I - kompilacja do formatu DVI}
	\begin{columns}[T]
		\begin{column}{.70\textwidth}
			\begin{itemize}
				\item DVI z ang. \emph{DeVice Independent}
				\item format binarny opisujący wizualny układ dokumentu
				\item odczyt i tworzenie łatwe do zaprogramowania
				\item interpretacja poleceń maszyny wirtualnej
			\end{itemize}
		\end{column}
		\begin{column}{.3\textwidth}
			\includegraphics[width=\textwidth]{synch3.png}
		\end{column}
	\end{columns}
\end{frame}

\begin{frame}
	\frametitle{Moduł prezentacji - wariant II - kompilacja do formatu PDF}
	\begin{columns}[T]
		\begin{column}{.70\textwidth}
			\begin{itemize}
				\item PDF z ang. \emph{Portable Document Format}
				\item opis dokumentu w okrojonym języku PostScript
				\item obsługa hiperłącz
				\item najpopularniejszy format dokumentów
				\item mechanizm SyncTeX
			\end{itemize}
		\end{column}
		\begin{column}{.3\textwidth}
			\includegraphics[width=\textwidth]{synch2.png}
		\end{column}
	\end{columns}
\end{frame}

\begin{frame}
	\frametitle{Moduł prezentacji - możliwość integracji aplikacji zewnętrznej}
	\begin{figure}[htp]
		\centering
		\includegraphics[width=.85\textwidth]{sumatra.png}
		\label{sum}
	\end{figure}
\end{frame}

\section{Dotychczasowe wyniki pracy}
\subsection{}

\begin{frame}
	\frametitle{SDI-1 czyli poprzedni semestr}
		PDI-1 poświęcona została:
		\begin{itemize}
			\item zapoznaniu z technologią \LaTeX\ 
			\item zapoznaniu z klasą dokumentów \beamer\footnote{Na podstawie \emph{T. Tantau, J. Wright, and V. Miletic. The beamer class user guide for version 3.36., 2013 } dostępne online na \url{http://www.ctan.org}}
			\item przygotowaniu wzorców dla szablonów slajdów
			\item typowaniu funkcjonalności \LaTeX-a, które należy przenieść do aplikacji
		\end{itemize}
\end{frame}

\begin{frame}
	\frametitle{SDI-2 czyli początek bieżącego semstru}
	Początek PDI-2 poświęcono na:
	\begin{itemize}
		\item prace koncepcyjne dotyczące architektury aplikacji
		\item uzupełnienie wiedzy dotyczącej \LaTeX-a m.in. mechanizm tabel i elementów blokowych
	\end{itemize}
\end{frame}

\section{W ramach podsumowania - plan prac}
\subsection{}

\begin{frame}
	\frametitle{Program minimum}
	\begin{enumerate}
		\item Wybór głównych narzędzi programistycznych
			\begin{itemize}
				\item język programowania
				\item biblioteka GUI
				\item inne narzędzia
			\end{itemize}
		\item Równoległa implementacja następujących obszarów:
			\begin{itemize}
				\item mechanizmy generowania kodu dokumentu klasy Beamer
				\item GUI umożliwiające wygodne tworzenie i edycję slajdów
			\end{itemize}
		\item Opracowanie i implementacja modułu prezentacji + integracja z etapem poprzednim
		\item Testowanie działania aplikacji
	\end{enumerate}
\end{frame}

\begin{frame}
	\frametitle{Kolejne cele}
	\begin{itemize}
		\item Mechanizm synchronizacji modułu prezentacji z GUI (SyncTeX)
		\item Mechanizm zapisu i odczytu plików .tex
		\item Przewidywana jest konieczność prac optymalizujących czas reakcji modułu prezentacji
	\end{itemize}	
\end{frame}

\begin{frame}[c]
	\Huge Dziękuję
\end{frame}

\end{document}
